%% kapa: exam=1 round=34 year=2023 period=1 form=A subject=real_estate subject_ko=부동산학원론 seq=23 number=103
%% source: Q-Net 감정평가사 기출문제 (cst00309) artlSeq=5211789 file=raw/1008/5211789/2023년도 제34회 감정평가사 1차시험 1교시.pdf
%% conversion: kordoc 4.24.51 -> tools/gen1.py (auto); review=proofread-sonnet; flags=none
\begin{question}{103}
부동산금융 및 투자에 관한 설명으로 옳지 않은 것은?
\begin{choices}
  \item 부동산금융은 부동산의 매입이나 매각, 개발 등과 관련하여 자금이나 신용을 조달하거나 제공하는 것을 말한다.
  \item 부동산의 특성과 관련하여 분할거래의 용이성과 생산의 장기성으로 인해 부동산금융은 부동산의 거래나 개발 등에서 중요한 역할을 하게 된다.
  \item 부동산투자에서 지분투자자가 대상 부동산에 가지는 권한을 지분권이라 하고, 저당투자자가 대상 부동산에 가지는 권한을 저당권이라 한다.
  \item 부동산보유자는 보유부동산의 증권화를 통해 유동성을 확보할 수 있다.
  \item 부동산금융이 일반금융과 다른 점으로는 담보기능과 감가상각 및 차입금 이자에 대한 세금감면이 있다.
\end{choices}
\end{question}
\answer{2}
